\section{Objetivo}

\paragraph{}
O objetivo deste teste de penetração é avaliar a segurança da aplicação web
disponibilizada no Laboratório de Segurança, identificando vulnerabilidades que
possam comprometer a confidencialidade, integridade e disponibilidade das
informações processadas pelo sistema.

\paragraph{}
A análise tem como base os resultados obtidos na atividade de varredura de
vulnerabilidades realizada com a ferramenta OWASP ZAP sobre o alvo
\texttt{http://testasp.vulnweb.com}. Durante a avaliação foram identificadas
falhas de alta severidade, incluindo \textit{SQL Injection}, \textit{Cross Site
Scripting (Reflected)}, \textit{Path Traversal} e \textit{External Redirect},
além de vulnerabilidades relacionadas à ausência de mecanismos de proteção como
Anti-CSRF, Content Security Policy (CSP) e proteção contra Clickjacking.

\paragraph{}
O teste buscará validar o impacto dessas vulnerabilidades, verificando a
possibilidade de acesso não autorizado a dados sensíveis, execução de código
malicioso no navegador dos usuários, manipulação de consultas ao banco de
dados, acesso indevido a arquivos internos do servidor e realização de ataques
de engenharia social por meio de redirecionamentos maliciosos.

\paragraph{}
Além da identificação e validação das vulnerabilidades encontradas, o teste tem
como finalidade avaliar a eficácia dos controles de segurança implementados na
aplicação e fornecer recomendações para mitigação dos riscos identificados,
contribuindo para o fortalecimento da postura de segurança do ambiente
analisado.
